\chapter[微分形式与外微分]{微分形式\\与外微分}
\label{ch:07-forms}
\ChapterPlan{周2.4（本地页码待核对）；梁上册第二版5.1及第5章习题2、5、6；同伦算子与局部引理另作补充}{正文及证明初稿；附8题解答与原创图，尚未独立数学审阅}

第6章的外代数运算只比较同一基点的向量。本章引入能比较邻近点变化的外微分，并说明它为什么不需要度规或联络。先在坐标中构造并证明无关选卡，再研究拉回、流与同伦；积分暂限于一维曲线和参数积分，流形上的一般积分与Stokes定理留到第8章。

\section{微分形式及局部表示}
\begin{definition}[微分形式与局部分量]
$M^n$上的光滑微分$k$形式（differential form）是$\Lambda^kT^*M$的光滑截面，其空间记为$\Omega^k(M)$。约定$\Omega^0(M)=C^\infty(M)$，$k<0$或$k>n$时为零空间。总空间$\Omega^\bullet(M)$按次数分级，乘法为逐点外积。
在卡$(U,x)$中，每个形式唯一写为
\[
 \omega=\sum_{i_1<\cdots<i_k}a_{i_1\cdots i_k}
 \dd x^{i_1}\wedge\cdots\wedge\dd x^{i_k}.
\]
系数为$\omega(\partial_{i_1},\ldots,\partial_{i_k})$。若对全部有序指标求和，则等价地写为
\[
 \omega=\frac1{k!}\omega_{i_1\cdots i_k}
 \dd x^{i_1}\wedge\cdots\wedge\dd x^{i_k},
 \qquad \omega_{i_1\cdots i_k}=\omega(\partial_{i_1},\ldots,\partial_{i_k}).
\]
\end{definition}
\begin{proposition}[形式的局部判别与换卡]\label{prop:forms-local}
形式光滑当且仅当上述递增指标系数光滑。若$y=y(x)$，则在新坐标中
\[
 \omega'_{a_1\cdots a_k}
 =\omega_{i_1\cdots i_k}
 \frac{\partial x^{i_1}}{\partial y^{a_1}}\cdots
 \frac{\partial x^{i_k}}{\partial y^{a_k}}.
\]
相容的局部形式可唯一拼成全局形式；外积与内插将光滑形式送到光滑形式。
\end{proposition}
\begin{proof}
由第6章外幂丛的构造，递增外积是局部光滑标架，截面判别即为各系数光滑。将$\partial_{y^a}=(\partial x^i/\partial y^a)\partial_{x^i}$逐槽代入，得到换卡式。全部有序指标式中的$k!$，来自每个递增指标组的$k!$个排列：分量与基外积的符号同时变化，乘积不变。
若局部形式在交上相同，每点的取值便与选块无关，定义出全局截面；光滑性局部可检验，因此拼接光滑且唯一。外积和内插的局部系数为有限和与乘积，故也光滑。
\end{proof}
\begin{example}[同一个面积形式的不同坐标]
在平面极坐标卡$r>0$上，$x=r\cos\theta,y=r\sin\theta$，于是
\[
 \dd x=\cos\theta\dd r-r\sin\theta\dd\theta,\qquad
 \dd y=\sin\theta\dd r+r\cos\theta\dd\theta.
\]
利用外积展开，$\dd x\wedge\dd y=r\dd r\wedge\dd\theta$。系数$r$并不说明几何形式改变，而是表示两个坐标标架的差别。极坐标仅在选定角区间的卡中使用；$\theta$不是穿孔平面上的全局实函数，$\dd\theta$的局部表达可相容，却未必是某个全局实函数的微分。
\end{example}
\begin{definition}[沿曲线的一形式积分]
设$\gamma:[a,b]\to M$分段光滑，$\alpha$为一形式，定义
\[
 \int_\gamma\alpha=\int_a^b\alpha_{\gamma(t)}(\dot\gamma(t))\dd t.
\]
每个光滑段分别积分后相加。这里使用单变量积分；在有限个坐标区间中计算的结果因逐点求值与选卡无关而一致。参数区间端点处只需相应单侧光滑延拓。
\end{definition}
\begin{proposition}[曲线积分与全微分]\label{prop:line-integral-exact}
保持方向的分段光滑微分同胚重参数化不改变积分，反向参数化使其变号。若$\alpha=\dd f$，则
\[
 \int_\gamma\dd f=f(\gamma(b))-f(\gamma(a)).
\]
因此恰当一形式沿任意闭曲线积分为零。
\end{proposition}
\begin{proof}
重参数化后速度乘参数导数，单变量换元公式给出两个方向的结论。对$\dd f$，链式法则将被积函数变为$(f\circ\gamma)'$，逐段用微积分基本定理，再将内部分段端点的函数值抵消，得到所列等式。这里没有调用尚未证明的Stokes定理。
\end{proof}
度规可把一形式与向量场联系起来，但本章的一形式、外积和曲线积分本来都不需要度规。局部系数与归一化练习见习题7.1。


\Needspace{18\baselineskip}
\section{外微分的构造与唯一性}
一形式$\alpha=P\dd x+Q\dd y$是否等于$\dd f$，首先受到混合偏导交换的限制$\partial_xQ=\partial_yP$。外微分将这一条件推广为任意次数的坐标无关运算。
\begin{theorem}[外微分的存在、性质与唯一性]\label{thm:exterior-derivative}
每个开集$U\subset M$上有与限制相容的实线性映射
\[
 \dd:\Omega^k(U)\to\Omega^{k+1}(U),
\]
满足：对函数$f$，$\dd f$为第4章定义的微分；$\dd^2=0$；且对$\alpha\in\Omega^p(U)$、$\beta\in\Omega^q(U)$，
\[
 \dd(\alpha\wedge\beta)=\dd\alpha\wedge\beta+(-1)^p\alpha\wedge\dd\beta.
\]
满足这些性质并与限制相容的运算族唯一。在坐标中，它由
\[
 \dd\omega=\sum_{I\text{递增}}\dd a_I\wedge\dd x^I,
 \qquad\dd a_I=\sum_j\partial_j a_I\dd x^j
\]
给出，其中$\dd x^I=\dd x^{i_1}\wedge\cdots\wedge\dd x^{i_k}$。
\end{theorem}
\begin{proof}
\textbf{在坐标块上定义并计算。}由形式的基展开唯一性，该公式定义了实线性运算，并把光滑系数变成光滑系数。对函数就是已知微分。计算$\dd^2\omega$时，只有
\[
 \sum_{I,j,\ell}\partial_\ell\partial_j a_I\,
 \dd x^\ell\wedge\dd x^j\wedge\dd x^I.
\]
$j=\ell$的项为零，$j\ne\ell$的两项由混合偏导相等和外积反对称性相消，所以平方为零。

对单项式$\alpha=a\dd x^I,\beta=b\dd x^J$，$\dd(ab)=b\dd a+a\dd b$。将$\dd b$移过前$p$个一次因子会产生$(-1)^p$，故得到带符号乘法式；当外积含重复指标时，两边同样按反对称性为零。按双线性展开即得一般式。

\textbf{证明换卡相容。}在同一交叠上取另一坐标$y$。用$x$坐标的运算记作$\dd_x$；已证$\dd_x^2=0$，故对每个坐标函数$y^a$有$\dd_x(\dd y^a)=0$。这里$\dd_x y^a=\dd y^a$由第4章的函数微分坐标无关性保证。将形式写成$\omega=\sum b_J\dd y^J$，用已证的乘法式得到
\[
 \dd_x\omega=\sum_J\dd b_J\wedge\dd y^J=\dd_y\omega.
\]
因此局部定义在交上一致，可唯一拼成全局形式，并自然与限制相容。全局的平方零和乘法式也由局部计算得到。

\textbf{唯一性。}若$D$是另一个满足条件的运算族，在坐标开集上$Dx^i=\dd x^i$，故$D(\dd x^i)=D^2x^i=0$。反复使用带符号乘法式，得到$D(a_I\dd x^I)=\dd a_I\wedge\dd x^I$。所以每块上$D=\dd$，由限制相容性得到全局一致。这个论证没有假设任意局部坐标函数都能作为全局实函数，也没有提前使用单位分解。
\end{proof}
外微分把函数的微分推广到所有形式，但它不是把每个张量分量随意求偏导。形式的交替性消去了换卡时的二阶项；第6章已经给出普通向量分量偏导不构成张量的反例。
\Needspace{18\baselineskip}
\begin{example}[低次公式与非坐标余标架]\label{ex:d-coframe}
在$\R^3$令$\alpha=P\dd x+Q\dd y+R\dd z$，并令
\[
 \beta=A\dd y\wedge\dd z+B\dd z\wedge\dd x+C\dd x\wedge\dd y.
\]
则
\begin{align*}
 \dd\alpha
 &=(\partial_xQ-\partial_yP)\dd x\wedge\dd y\\
 &\quad +(\partial_xR-\partial_zP)\dd x\wedge\dd z\\
 &\quad +(\partial_yR-\partial_zQ)\dd y\wedge\dd z,\\
 \dd\beta
 &=(\partial_xA+\partial_yB+\partial_zC)\dd x\wedge\dd y\wedge\dd z.
\end{align*}
这些公式先是形式等式，未借度规把形式与向量场等同。
若取平面上的余标架$\theta^1=\dd x,\theta^2=e^x\dd y$，则
\[
 \dd\theta^2=e^x\dd x\wedge\dd y=\theta^1\wedge\theta^2\ne0.
\]
因此对$\omega=a\theta^2$，有$\dd\omega=\dd a\wedge\theta^2+a\dd\theta^2$，不能漏掉第二项。只有坐标微分$\dd x^i$保证外微分为零。
\end{example}
\begin{definition}[闭形式与恰当形式]
若$\dd\omega=0$，称$\omega$闭（closed）；若$\omega=\dd\eta$，称恰当（exact），也称正合或全微分形式。由$\dd^2=0$，恰当形式必闭。零形式闭等价于局部常值；在连通流形上才进一步为常数。按$\Omega^{-1}=0$约定，恰当零形式只有零函数。
\end{definition}
闭零形式局部常值，是因为在凸小坐标球内沿线段积分$\dd f=0$便得函数值不变；连通空间中局部常值函数的各个值域原像开闭，所以只能取一个值。正次数下“闭是否推出恰当”将在本章最后回答。
\Needspace{12\baselineskip}
\begin{proposition}[外微分的无坐标求值式]\label{prop:koszul-d}
对$\omega\in\Omega^k(M)$及光滑场$X_0,\ldots,X_k$，有
\begin{align*}
 (\dd\omega)(X_0,\ldots,X_k)
 &=\sum_{i=0}^k(-1)^i X_i\bigl(\omega(X_0,\ldots,\widehat X_i,\ldots,X_k)\bigr)\\
 &\quad+\sum_{0\leq i<j\leq k}(-1)^{i+j}
 \omega([X_i,X_j],X_0,\ldots,\widehat X_i,\ldots,\widehat X_j,\ldots,X_k).
\end{align*}
第二行的括号放在第一个槽，其他场保持原次序。$k=1$时即
\[
 \dd\alpha(X,Y)=X(\alpha(Y))-Y(\alpha(X))-\alpha([X,Y]).
\]
\end{proposition}
\begin{proof}
先检查右边对各$X_j$是函数线性的。将$X_j$换成$fX_j$，第一和式中$i\ne j$的求导会多出$X_i(f)$乘删去$X_i$后的形式值；第二和式用$[X_i,fX_j]=f[X_i,X_j]+X_i(f)X_j$及$[fX_j,X_i]=f[X_j,X_i]-X_i(f)X_j$展开。把插入的$X_j$从第一个槽移回它原来的位置，所需交换的符号使这些额外项逐一抵消。对$i=j$的项只有整体因子$f$。因此右边为$C^\infty$多重线性。交换两个相邻输入时，第一和式按删去位置重新编号，第二和式同时使用括号反对称性，整体变号，故它确实是交替形式的求值。
现在只需在一组局部坐标基场上核对。坐标场的括号为零，第一和式就是分量系数的交替偏导，等于定理\ref{thm:exterior-derivative}的坐标定义。两侧均函数多重线性，故对所有场成立；$k=0$时退化为函数微分的定义。习题7.2用非坐标场检查括号项的必要性。
\end{proof}


\Needspace{15\baselineskip}
\section{拉回与Cartan公式}
\begin{definition}[形式的拉回]
若$F:M\to N$光滑，对$\omega\in\Omega^k(N)$定义
\[
 (F^*\omega)_p(v_1,\ldots,v_k)
 =\omega_{F(p)}(\dd F_pv_1,\ldots,\dd F_pv_k).
\]
零次时$F^*f=f\circ F$。它与第6章的协变张量拉回相同，不要求$F$单射、满射或局部可逆。
\end{definition}
\begin{theorem}[拉回与外微分的自然性]\label{thm:d-pullback}
拉回保持次数、外积、恒等与反向复合，且
\[
 \dd(F^*\omega)=F^*(\dd\omega).
\]
因此闭形式与恰当形式均拉回为相应类型的形式。
\end{theorem}
\begin{proof}
在局部坐标中$F^*(a_I\dd y^I)=(a_I\circ F)\dd F^{i_1}\wedge\cdots\wedge\dd F^{i_k}$，其中$F^a=y^a\circ F$。这些系数光滑，故拉回光滑。逐点的线性代数给出外积相容，链式法则给出反向复合。
对上述局部项求外微分，乘法式中$\dd(\dd F^a)=0$，只剩
\[
 \dd(a_I\circ F)\wedge\dd F^{i_1}\wedge\cdots\wedge\dd F^{i_k}
 =F^*(\dd a_I\wedge\dd y^I),
\]
用了函数微分的链式法则。按项相加即得自然性。目标任意点附近可取卡，利用$F$连续性在源上缩小，故这些局部计算覆盖全部$M$。最后若$\dd\omega=0$，自然性给$\dd F^*\omega=0$；若$\omega=\dd\eta$，则$F^*\omega=\dd F^*\eta$。
\end{proof}
\begin{figure}[htbp]
\centering\DGDiagram{forms-pullback}

\caption{原创交换图：任意光滑映射$F:M\to N$使拉回与外微分交换。形式的拉回方向与$F$相反；图中没有假设$F$可逆，也没有用图表代替自然性的证明。}
\label{fig:forms-pullback}
\end{figure}
对光滑向量场$X$，第6章逐点定义的内插$\iota_X:\Omega^k\to\Omega^{k-1}$是次数$-1$的带符号导子。它使形式少一个输入槽，却不求导系数。另一种运算是沿$X$的流比较形式。
\begin{definition}[形式的Lie导数]
设$\Phi_t$为$X$的极大局部流，定义
\[
 \mathcal L_X\omega=\left.\frac{\dd}{\dd t}\right|_{t=0}\Phi_t^*\omega.
\]
每点附近均有统一短时间的流，故可在固定开邻域中对光滑系数求导。所得形式局部相容，不需要$X$完备。拉回将$\Phi_t(p)$的形式放到固定的$T_pM$上，因此比较发生在同一纤维中。
\end{definition}
\begin{lemma}[Lie导数的局部生成规则]\label{lem:lie-forms}
$\mathcal L_X$是次数零的实线性导子，满足
\[
 \mathcal L_Xf=Xf,\qquad\mathcal L_X(\dd f)=\dd(Xf),
\]
\[
 \mathcal L_X(\alpha\wedge\beta)
 =(\mathcal L_X\alpha)\wedge\beta+\alpha\wedge\mathcal L_X\beta.
\]
\end{lemma}
\begin{proof}
拉回保持外积，对$\Phi_t^*(\alpha\wedge\beta)=\Phi_t^*\alpha\wedge\Phi_t^*\beta$在零处求导即得乘法式。函数上由流的速度为$X$得到第一式。由自然性，$\Phi_t^*\dd f=\dd(f\circ\Phi_t)$，对$t$求导后可交换光滑混合偏导，得到第二式。实线性与次数保持由定义直接得到。局部形式由函数与$\dd x^i$生成，因此这些规则确定$\mathcal L_X$在所有形式上的值。
\end{proof}
\Needspace{12\baselineskip}
\begin{theorem}[Cartan公式]\label{thm:cartan}
对任意光滑场$X$和微分形式$\omega$，
\[
 \mathcal L_X\omega=\iota_X\dd\omega+\dd\iota_X\omega.
\]
对函数约定$\iota_Xf=0$，故公式也适用于零次形式。
\end{theorem}
\begin{proof}
记右边算子为$D=\iota_X\dd+\dd\iota_X$。将外微分的带符号乘法式及内插的带符号乘法式依次代入$D(\alpha\wedge\beta)$。含$\dd\alpha\wedge\iota_X\beta$的交叉项符号分别为$(-1)^{p+1}$与$(-1)^p$，相消；含$\iota_X\alpha\wedge\dd\beta$的两项符号为$(-1)^p$与$(-1)^{p-1}$，也相消。其余项恰为
\[
 D\alpha\wedge\beta+\alpha\wedge D\beta.
\]
所以$D$是次数零导子。在函数上$Df=\iota_X\dd f=Xf$；在$\dd f$上由$\dd^2=0$得$D(\dd f)=\dd(Xf)$。它与引理\ref{lem:lie-forms}的$\mathcal L_X$在局部生成元$f,\dd x^i$上相同，故通过乘法规则对所有局部形式相同，继而全局相同。
\end{proof}
\begin{proposition}[求值式与Cartan交换关系]\label{prop:cartan-identities}
对光滑场$Y_1,\ldots,Y_k$，
\begin{align*}
 (\mathcal L_X\omega)(Y_1,\ldots,Y_k)
 &=X\bigl(\omega(Y_1,\ldots,Y_k)\bigr)\\
 &\quad-\sum_{j=1}^k\omega(Y_1,\ldots,[X,Y_j],\ldots,Y_k).
\end{align*}
此外有
\[
 \mathcal L_X\dd=\dd\mathcal L_X,\qquad
 \mathcal L_X\iota_Y-\iota_Y\mathcal L_X=\iota_{[X,Y]},
\]
\[
 \mathcal L_X\mathcal L_Y-\mathcal L_Y\mathcal L_X=\mathcal L_{[X,Y]}.
\]
\end{proposition}
\begin{proof}
求值式右边在$Y_j$换成$fY_j$时，第一项多出$X(f)$乘形式值，括号项由$[X,fY_j]=X(f)Y_j+f[X,Y_j]$产生相反项，故函数多重线性；它也交替。在外积上展开求值，按槽位分组可见它满足次数零导子的乘法规则。函数上等于$Xf$，对$\dd f$及一个场$Y$，其值为$X(Yf)-[X,Y]f=Y(Xf)$，即$\dd(Xf)(Y)$。因此它与$\mathcal L_X$在局部生成元上一致，得到求值式。这也与第5章沿流拉回场的括号符号相容。

第一交换关系直接由Cartan公式和$\dd^2=0$得到，两边都等于$\dd\iota_X\dd$。第二关系的左边是次数$-1$导子，这可用两种乘法式展开并消去交叉项验证；它在函数上为零，在$\dd f$上为$X(Yf)-Y(Xf)=[X,Y]f$，恰与$\iota_{[X,Y]}$相同。第三关系两边为次数零导子，在函数上都是$[X,Y]f$；在$\dd f$上，利用第一关系，两边都是$\dd([X,Y]f)$。再用局部生成元判别，得到两式对所有形式成立。
\end{proof}
\begin{example}[沿伸缩流的次数权重]
在$\R^n$令$R=x^i\partial_i$，流为$\Phi_t(x)=e^tx$。因$\Phi_t^*\dd x^i=e^t\dd x^i$，常系数$k$形式满足$\mathcal L_R\omega=k\omega$。若系数$a(x)$是$m$次齐次函数，则$\mathcal L_R(a\dd x^I)=(m+k)a\dd x^I$。不能只对系数求方向导数而忽略基形式本身的变化。习题7.3、7.8检查Cartan公式中的两项与交换关系。
\end{example}


\Needspace{15\baselineskip}
\section{Poincare引理与同伦算子}
闭形式的原函数不仅需要存在性论证，还希望有可以计算的构造。先建立一个比较两端拉回的恒等式，再把它用于把星形域收缩到一点的同伦。
\begin{theorem}[同伦算子公式]\label{thm:homotopy-operator}
设$H(t,p)$在$[0,1]\times M$的某个开邻域上光滑并取值于$N$，记$H_t(p)=H(t,p)$、$j_t(p)=(t,p)$。定义次数$-1$的线性算子
\[
 K_H\omega=\int_0^1 j_t^*(\iota_{\partial_t}H^*\omega)\dd t,
 \qquad K_Hf=0.
\]
积分逐点在固定的形式纤维中进行。则
\[
 \dd K_H\omega+K_H\dd\omega=H_1^*\omega-H_0^*\omega.
\]
\end{theorem}
\begin{proof}
在时间与$M$的乘积坐标中，唯一分解
\[
 H^*\omega=\dd t\wedge\alpha_t+\beta_t,
 \qquad \alpha_t=j_t^*(\iota_{\partial_t}H^*\omega),\quad\beta_t=H_t^*\omega.
\]
每点$p$附近可用同一组$M$坐标写出这些系数。对紧时间段$[0,1]$，开邻域条件允许缩小$p$的邻域，使这些系数在一个稍大的时间段内光滑；由参数积分的微分法则，$K_H\omega$光滑，且$\dd_M$可穿过该积分。
由外微分的乘法式，
\[
 \dd(H^*\omega)
 =\dd t\wedge(\partial_t\beta_t-\dd_M\alpha_t)+\dd_M\beta_t.
\]
自然性使左边等于$H^*\dd\omega$。内插$\partial_t$后限制到时间片并积分，得到
\[
 K_H\dd\omega=\int_0^1(\partial_t\beta_t-\dd_M\alpha_t)\dd t
 =\beta_1-\beta_0-\dd_MK_H\omega.
\]
这就是所需公式。若$\omega=f$为零形式，分解只有$\beta_t$，对$\dd f$应用同样计算，仍给出$K_H\dd f=H_1^*f-H_0^*f$，没有负次数形式问题。
\end{proof}
这里不要求存在一个对所有$p$统一的时间延长区间；只需在$[0,1]\times M$周围有开邻域，证明逐点使用紧时间段即可。这一细节适用于任意开星形域，而不只是整个欧氏空间。
\begin{theorem}[星形域上的Poincare引理]\label{thm:poincare}
设$U\subset\R^n$开，关于$a\in U$星形，即$a+t(x-a)\in U$对$x\in U$、$0\leq t\leq1$成立。则每个正次数闭形式$\omega\in\Omega^k(U)$均恰当。可取原形式$K\omega$，其中
\[
 (K\omega)_x(v_1,\ldots,v_{k-1})
 =\int_0^1t^{k-1}\omega_{a+t(x-a)}(x-a,v_1,\ldots,v_{k-1})\dd t.
\]
任意流形上的正次数闭形式都局部恰当。零次结论是闭函数局部常值，而不是任意闭函数都恰当。
\end{theorem}
\begin{proof}
令$H(t,x)=a+t(x-a)$。集合$\{(t,x):x\in U,H(t,x)\in U\}$是包含$[0,1]\times U$的开集，因此满足上一定理的延拓条件。时间方向的微分为$x-a$，每个空间输入的微分为$tv_j$；逐槽线性正好给出因子$t^{k-1}$。$t=0$处$k=1$时因子为一，其余正次数时为非负整数次幂，积分光滑，并无$1/t$奇性。
此时$H_1=\operatorname{id}_U$，$H_0$是常值映射。正次数的形式在常值映射下拉回为零，因为其切映射为零；同伦公式于是成为$\dd K\omega+K\dd\omega=\omega$。当$\dd\omega=0$，得到$\omega=\dd K\omega$。对零形式，则$K\dd f=f-f(a)$。
一般流形每点有坐标像为欧氏小球的卡。把闭形式移到坐标球，用上式得到局部原形式，再拉回流形卡域即可。局部原形式未保证能在不同卡上无差别拼接为一个全局原形式。
\end{proof}
\begin{figure}[htbp]
\centering\DGDiagram{radial-homotopy}

\caption{原创示意：星形域内的径向同伦$H(t,x)=tx$（取中心$a=0$）。时间方向微分为$x$，每个空间输入变成$tv$，从而在$k$形式的同伦算子中产生$t^{k-1}$。图画出固定$x$的一条收缩路径；整个路径须留在定义域中。}
\label{fig:radial-homotopy}
\end{figure}
\begin{example}[常系数形式的显式原形式]
在$\R^n$以零点为中心，对常系数$k$形式$\omega$、$k\geq1$，有$K\omega=\frac1k\iota_R\omega$，其中$R_x=x$。因为积分只剩$\int_0^1t^{k-1}\dd t=1/k$，所以例如
\[
 \dd x\wedge\dd y=\dd\!\left(\frac{x\dd y-y\dd x}{2}\right).
\]
一般一形式$\alpha=\alpha_i\dd x^i$若闭，其星形域原函数可取$f(x)=\int_0^1\alpha_i(tx)x^i\dd t$；这是一个可计算公式，而不仅是存在性陈述。非闭形式则必须保留同伦公式中的$K\dd\alpha$项，见习题7.5。
\end{example}
\Needspace{15\baselineskip}
\begin{example}[穿孔平面上的闭而不恰当形式]\label{ex:angular-form}
在$M=\R^2\setminus\{0\}$取
\[
 \eta=\frac{-y\dd x+x\dd y}{x^2+y^2}.
\]
若写$\eta=P\dd x+Q\dd y$，直接求导得
\[
 \partial_xQ=\frac{y^2-x^2}{(x^2+y^2)^2}=\partial_yP,
\]
故$\dd\eta=0$。但在逆时针单位圆$\gamma(t)=(\cos t,\sin t)$、$0\leq t\leq2\pi$上，$\gamma^*\eta=\dd t$，所以$\int_\gamma\eta=2\pi$。若$\eta=\dd f$，命题\ref{prop:line-integral-exact}却要求此闭曲线积分为零，矛盾。因而$\eta$只局部为角函数的微分。
径向收缩到原点不适用于这个域，因为原点被删去；将星形域结论直接套用到穿孔平面，违反了同伦始终取值于定义域的假设。
\end{example}
\begin{figure}[htbp]
\centering\DGDiagram{angular-obstruction}

\caption{原创示意：穿孔平面上的角形式沿逆时针圆周有积分$2\pi$，在局部角卡中却等于$\dd\theta$。橙色空圈表示被删去的原点，圆周箭头给出参数方向；图不把局部角函数画成全局单值实函数。}
\label{fig:angular-obstruction}
\end{figure}
同伦公式会在第9章用于de Rham上同调，但本章不预先使用上同调、Stokes或拓扑同伦的光滑逼近定理来证明它。所有收缩与同伦均以明确的光滑映射及其定义域为条件。


\Needspace{15\baselineskip}
\section{微分形式习题及详细解答}
习题7.1、7.2、7.3分别改述或扩展梁书上册第二版第5章习题2、5、6，原题均在印刷页130（PDF物理页146）。其余5题为原创补充。本文沿用第6章的完整行列式外积约定；其与梁书5.1的阶乘归一化已核对。原题号与新增部分见\texttt{references/EXERCISES.md}，全部解答由本工程编写。

\Needspace{12\baselineskip}
\subsection*{习题7.1：四维形式的分量展开}
改述梁书第5章习题2：设$V$为四维实空间，基为$e_1,\ldots,e_4$，对偶基为$e^1,\ldots,e^4$。分别写出一般一形式、三形式与四形式的基展开，说明系数的求值含义，并核对全指标求和时的阶乘。
\begin{proof}[解答]
一形式为$\alpha=\sum_{i=1}^4a_i e^i$，其中$a_i=\alpha(e_i)$。三形式为
\begin{align*}
 \beta&=b_{123}e^1\wedge e^2\wedge e^3+b_{124}e^1\wedge e^2\wedge e^4\\
 &\quad+b_{134}e^1\wedge e^3\wedge e^4+b_{234}e^2\wedge e^3\wedge e^4,
\end{align*}
其中$b_{ijk}=\beta(e_i,e_j,e_k)$对$i<j<k$成立；其余有序分量按排列变号，重复指标时为零。四形式为$\gamma=c\,e^1\wedge e^2\wedge e^3\wedge e^4$，$c=\gamma(e_1,e_2,e_3,e_4)$。
若对全部有序指标求和，同一递增组的$k!$项因系数和基外积同时变号而相同，故必须分别写$\beta=\frac1{3!}b_{ijk}e^i\wedge e^j\wedge e^k$与$\gamma=\frac1{4!}\gamma_{ijkl}e^i\wedge e^j\wedge e^k\wedge e^l$。递增求和时再加这些系数，就会重复除以阶乘。
\end{proof}

\Needspace{12\baselineskip}
\subsection*{习题7.2：非坐标场中的括号项}
以梁书第5章习题5为出发点，证明一形式的求值恒等式
$\dd\alpha(X,Y)=X(\alpha(Y))-Y(\alpha(X))-\alpha([X,Y])$。
再取$X=\partial_x,Y=e^x\partial_y,\alpha=e^{-x}\dd y$，逐项计算，判断能否省掉括号项。
\begin{proof}[解答]
写$\alpha=a_i\dd x^i$、$X=X^i\partial_i,Y=Y^j\partial_j$。展开两个方向导数后相减，包含$\partial_iY^j$与$\partial_jX^i$的项组成$\alpha([X,Y])$；余下为
\[
 X^iY^j(\partial_i a_j-\partial_j a_i)=\dd\alpha(X,Y).
\]
移项即得恒等式，证明未要求$X,Y$为坐标场。
在给定例中$\alpha(Y)=1,\alpha(X)=0$，所以前两个导数都为零。$[X,Y]=Y$，故括号项为$-1$。另一方面，$\dd\alpha=-e^{-x}\dd x\wedge\dd y$，在$(X,Y)$上也取值$-1$。若删去括号项，就错误得到零。坐标场的特殊简化不能当成任意标架下的公式。
\end{proof}

\Needspace{12\baselineskip}
\subsection*{习题7.3：Cartan公式的两项}
改述并扩展梁书第5章习题6：从$\mathcal L_X=\iota_X\dd+\dd\iota_X$证明$\mathcal L_X\dd=\dd\mathcal L_X$。在平面取$X=x\partial_x+y\partial_y$、$\alpha=x\dd y-y\dd x$，计算$\iota_X\alpha$、$\dd\alpha$、$\mathcal L_X\alpha$及$\mathcal L_X\dd\alpha$，并与流的拉回核对。
\begin{proof}[解答]
由$\dd^2=0$，$\mathcal L_X\dd=\dd\iota_X\dd=\dd\mathcal L_X$。本例$\iota_X\alpha=xy-yx=0$，而$\dd\alpha=2\dd x\wedge\dd y$，所以
\[
 \mathcal L_X\alpha=2(x\dd y-y\dd x)=2\alpha,
 \qquad \mathcal L_X\dd\alpha=\dd(2\alpha)=4\dd x\wedge\dd y.
\]
$X$的流为$\Phi_t(x,y)=(e^tx,e^ty)$。拉回时系数与一次基形式分别乘$e^t$，故$\Phi_t^*\alpha=e^{2t}\alpha$，零时导数为$2\alpha$；$\dd\alpha$的两个一次因子也使拉回乘$e^{2t}$，导数为$4\dd x\wedge\dd y$。若只对系数作方向导数，会少掉基形式变化的贡献。
\end{proof}

\Needspace{12\baselineskip}
\subsection*{习题7.4：退化映射下的自然性}
设$F:\R^2\to\R^2$，$F(s,t)=(x,y)=(s^2,t)$，目标一形式$\beta=x\dd y$。求$F^*\beta$及$F^*\dd\beta$，核对自然性，并解释$s=0$处切映射退化为何不妨碍拉回。
\begin{proof}[解答]
直接有$F^*\beta=s^2\dd t$，从而$\dd F^*\beta=2s\dd s\wedge\dd t$。目标$\dd\beta=\dd x\wedge\dd y$，而$F^*\dd x=2s\dd s,F^*\dd y=\dd t$，故$F^*\dd\beta$也等于$2s\dd s\wedge\dd t$。
在$s=0$，矩阵$DF=\operatorname{diag}(0,1)$秩一，二形式拉回在那里为零。拉回只需把源切向量送入目标形式，并不需要逆切映射；退化会使部分输入信息丢失，却不会使定义失效。本例也不满足单射条件，但两点同像不影响各自基点上的拉回定义。
\end{proof}

\Needspace{12\baselineskip}
\subsection*{习题7.5：非闭形式的同伦余项}
在$\R^2$使用以零点为中心的径向同伦。对$\omega=x\dd y$计算$K\omega$、$\dd K\omega$、$K\dd\omega$，核对$\dd K\omega+K\dd\omega=\omega$。为何不能无条件把$K\omega$称作原函数？
\begin{proof}[解答]
一形式上$K\omega(x,y)=\int_0^1(tx)y\dd t=xy/2$，所以$\dd K\omega=(y\dd x+x\dd y)/2$。又$\dd\omega=\dd x\wedge\dd y$，同伦公式给
\[
 K\dd\omega=\int_0^1t(x\dd y-y\dd x)\dd t
 =\frac{x\dd y-y\dd x}{2}.
\]
两者相加为$x\dd y$，确实等于$\omega$。但$\dd\omega\ne0$，余项不能删掉；因此$\dd K\omega\ne\omega$，$K\omega$不是它的原函数。Poincare引理的闭性假设是这个构造能够给出原形式的关键。
\end{proof}

\Needspace{12\baselineskip}
\subsection*{习题7.6：角形式的局部原函数与周期}
对例\ref{ex:angular-form}中的$\eta$，在右半平面写出一个原函数，说明它为何不能延伸为穿孔平面上满足$\dd f=\eta$的单值光滑实函数；对半径$R>0$的逆时针圆再计算积分。
\begin{proof}[解答]
右半平面$x>0$可取$f(x,y)=\arctan(y/x)$。直接微分得$\dd f=(-y\dd x+x\dd y)/(x^2+y^2)=\eta$。在其他角卡中可用相应角分支，交叠的连通部分上原函数相差常数，但绕原点一周会积累$2\pi$，不能拼成全局单值实函数。
更直接地，令$\gamma_R(t)=(R\cos t,R\sin t)$。代入后分子为$R^2\dd t$，分母为$R^2$，所以$\int_{\gamma_R}\eta=2\pi$，与$R$无关。若全局$\dd f=\eta$，曲线积分应为端点函数值之差，即零，矛盾。这个判断只用一维积分，并未把尚未建立的Stokes定理用于穿孔域。
\end{proof}

\Needspace{12\baselineskip}
\subsection*{习题7.7：平移同伦的端点差}
设$H(t,x,y)=(x+t,y)$，$0\leq t\leq1$，目标形式$\omega=x\dd y$。用$K_H$的定义分别计算$K_H\omega$、$K_H\dd\omega$及$H_1^*\omega-H_0^*\omega$。
\begin{proof}[解答]
在三变量$(t,x,y)$上，$H^*\omega=(x+t)\dd y$，不含$\dd t$，所以$K_H\omega=0$。但
\[
 H^*\dd\omega=(\dd x+\dd t)\wedge\dd y,
 \qquad\iota_{\partial_t}H^*\dd\omega=\dd y.
\]
限制到时间片后积分，得到$K_H\dd\omega=\dd y$。端点差为$(x+1)\dd y-x\dd y=\dd y$，因此正好与同伦公式的左边相同。一般同伦的算子来自时间方向内插；不能把径向公式中的$t^{k-1}$直接挪用到任意同伦。
\end{proof}

\Needspace{12\baselineskip}
\subsection*{习题7.8：Cartan交换关系的符号}
在平面取$X=\partial_x,Y=x\partial_y$。用$\alpha=\dd y$核对$[\mathcal L_X,\iota_Y]\alpha=\iota_{[X,Y]}\alpha$，再在$f=xy$上核对$[\mathcal L_X,\mathcal L_Y]f=\mathcal L_{[X,Y]}f$。
\begin{proof}[解答]
$[X,Y]=\partial_y$。$\iota_Y\alpha=x$，所以$\mathcal L_X\iota_Y\alpha=1$；$\mathcal L_X\alpha=\dd(Xy)=0$，故$\iota_Y\mathcal L_X\alpha=0$。左边为$1$，右边$\iota_{\partial_y}\dd y$也为$1$。
对$f=xy$，先有$Yf=x^2$、$Xf=y$。于是$X(Yf)=2x$、$Y(Xf)=x$，两者之差为$x$；右边$\partial_y(xy)=x$，再次相同。交换算子顺序会改号，因此应始终与$[X,Y]=XY-YX$的约定配合。
\end{proof}

